
\section{Introduction}
\label{intro}

Modern artificial intelligence (AI) workloads rely on specialized hardware accelerators to achieve high performance and energy efficiency. 
field-programmable gate arrays (FPGAs) have emerged as an attractive option for this purpose due to their reconfigurability and ability to provide customized computation pipelines tailored to specific AI workloads. 
However, designing efficient FPGA-based accelerators for AI workloads remains a challenging and time-consuming process that requires extensive expertise in hardware-software co-design, high-level synthesis (HLS), and system integration.


% The design of AI accelerators involves exploring a large hardware design space that includes architectural parameters, data-flow strategies, and memory hierarchies.
AI accelerators comprise architectural parameters, data-flow strategies, and memory hierarchies; together, these factors create a large hardware design space that must be explored during the design process.
This exploration is often performed manually or supported by a limited number of automated tools~\cite{cesaretti2025rapid, tiwari2025hardware}, leading to prolonged development cycles and effort.
Furthermore, each design iteration may require costly HLS and evaluation stages, slowing down the exploration of the hardware-software co-design space~\cite{gibsonDLASConceptualModel2025}.

The SECDA (SystemC Enabled Co-Design of DNN Accelerators) methodology~\cite{haris2021secda} was proposed to reduce the development time and effort of FPGA-based accelerators by enabling hardware-software co-design through SystemC simulation and hardware synthesis workflows. 
This approach allows developers to rapidly iterate on accelerator designs while maintaining a unified development flow.


Building upon the SECDA TFLite, the SECDA-DSE framework aims to further automate the accelerator design process by incorporating design space exploration (DSE) techniques through the effective use of large language models (LLMs).
SECDA-DSE leverages the reasoning capabilities of LLMs to guide the exploration of hardware design parameters, generate candidate accelerator architectures, and iteratively refine designs based on evaluation and performance feedback. 
We discuss the framework architecture, including the two major component, the DSE Explorer and the LLM Stack, and outline how these components enable iterative exploration, evaluation, and refinement of new designs. 
We present the vision and early design of SECDA-DSE along with preliminary results through verification of a generated accelerator design through HLS on a Zynq-7000 FPGA.


%Finally, the paper discuss the potential future extending of SECDA-DSE toward a fully automated hardware design exploration pipeline.